\chapter{Équations, inéquations et signes}
\section*{Exercice 6.1 -- Premier degré}
Résoudre $5x-7=2x+8$, puis $3(x-1)=3x+2$.

\section*{Exercice 6.2 -- Produit nul}
Résoudre $(2x-1)(x+4)=0$ et $x^2=12$.

\section*{Exercice 6.3 -- Quotient}
Résoudre $(x+3)/(x-1)=2$, puis $(x^2-1)/(x-1)=2$.

\section*{Exercice 6.4 -- Signes}
Résoudre $(x-3)(x+2)>0$ et $(x-3)/(x+2)\leq0$.

\section*{Exercice 6.5 -- Budget}
Un forfait coûte $14$ euros puis $3$ euros par séance. Avec un budget de $50$ euros, combien de séances au maximum peut-on payer ?

\section*{Exercice 6.6 -- Équivalence ?}
Un élève résout $x(x-2)=0$ en divisant par $x$ et conclut $x=2$. Expliquer et corriger son erreur.
